\documentclass[a4paper,11pt]{article}
\usepackage[T1]{fontenc}
\usepackage[utf8]{inputenc}
\usepackage[italian]{babel}
\usepackage{amsmath,amssymb}
\usepackage{booktabs}
\usepackage{longtable}
\usepackage{geometry}
\usepackage{hyperref}
\usepackage{xcolor}
\usepackage{enumitem}
\usepackage{graphicx}

\geometry{margin=2.5cm}

\newcommand{\desc}[1]{\subsection{#1}}
\newcommand{\formula}{\paragraph{Formula.}}
\newcommand{\ascolta}{\paragraph{Cosa ascolta.}}
\newcommand{\range}{\paragraph{Valori tipici.}}
\newcommand{\rif}{\paragraph{Riferimenti.}}
\newcommand{\oss}{\paragraph{Osservazioni dai test.}}

\title{Descrittori spettrali per l'analisi del timbro\\[0.5em]
\large Dispensa di studio per \emph{interfantasia}}
\author{Francesco Ferrante}
\date{Aprile 2026}

\begin{document}
\maketitle

\tableofcontents
\newpage

% ============================================================
\section{Introduzione}
% ============================================================

Questa dispensa raccoglie i descrittori spettrali studiati e implementati per
il sistema \emph{interfantasia}. L'obiettivo è comprendere cosa ciascun
descrittore cattura del segnale audio, per poi decidere come mappare i
descrittori verso i parametri del filtro bicomb.

Tutti i descrittori vengono calcolati frame per frame su una
rappresentazione FFT del segnale. I parametri di analisi sono:
\begin{itemize}[nosep]
  \item FFT size: 8192 campioni
  \item Sample rate: 96\,000 Hz (risoluzione in frequenza: $\sim$11.7 Hz per bin)
  \item Finestra: Hann, overlap 50\%
  \item Banda analizzata: 0--10\,000 Hz
\end{itemize}

I segnali di test sono generati in Csound e comprendono: sinusoidi pure,
noise bianco, sinusoidi in \texttt{tanh} (distorsione a drive variabile),
noise filtrato bandpass, sintesi FM, miscele sinusoide/noise, glissandi
e coppie di sinusoidi.

I descrittori sono organizzati in quattro categorie:
\begin{enumerate}[nosep]
  \item \textbf{Forma dello spettro}: descrivono dove si concentra l'energia
  \item \textbf{Distribuzione}: descrivono come l'energia è distribuita
  \item \textbf{Tonalità}: misurano quanto il segnale è periodico o rumoroso
  \item \textbf{Dinamica}: misurano il cambiamento nel tempo
\end{enumerate}

\subsection{Notazione}

In tutte le formule:
\begin{itemize}[nosep]
  \item $X(k)$ è l'ampiezza (magnitudine) della $k$-esima bin della FFT
  \item $f(k)$ è la frequenza corrispondente alla bin $k$
  \item $N$ è il numero di bin considerate
  \item $\mu_1$ è il centroide spettrale (momento primo)
  \item $\sigma$ è lo spread spettrale (deviazione standard)
\end{itemize}


% ============================================================
\section{Forma dello spettro}
% ============================================================

Questi descrittori indicano \emph{dove} si trova l'energia nello spettro.

% ---
\desc{Centroide spettrale (\emph{spectral centroid})}

\formula
\[
  \mu_1 = \frac{\displaystyle\sum_{k=0}^{N-1} f(k) \cdot X(k)}
               {\displaystyle\sum_{k=0}^{N-1} X(k)}
\]

\ascolta
Il ``baricentro'' dello spettro, la frequenza media pesata per ampiezza.
Corrisponde alla percezione di \textbf{brillantezza}: un centroide alto
indica un suono brillante, uno basso indica un suono cupo. È la prima
delle tre dimensioni percettive del timbro individuate da
Krimphoff et al.\ (1994).

\range
\begin{itemize}[nosep]
  \item Sinusoide 440 Hz: 440 Hz (coincide con la fondamentale)
  \item Noise bianco: 5\,000 Hz (centro della banda 0--10\,000 Hz)
  \item Tanh drive 20: 1\,876 Hz (le armoniche alzano il baricentro)
  \item FM indice 10: 4\,283 Hz (bande laterali estese)
\end{itemize}

\rif Peeters (2003, sez.~6.1.1), Lerch (2023, sez.~3.5.1), Krimphoff et al.\ (1994).

% ---
\desc{Spread spettrale (\emph{spectral spread})}

\formula
\[
  \sigma = \sqrt{\frac{\displaystyle\sum_{k=0}^{N-1} \big(f(k) - \mu_1\big)^2 \cdot X(k)}
                      {\displaystyle\sum_{k=0}^{N-1} X(k)}}
\]

\ascolta
La \textbf{larghezza di banda} dello spettro: quanto l'energia è
dispersa attorno al centroide. Un valore alto indica uno spettro ampio
(noise, molte armoniche), un valore basso indica energia concentrata
(sinusoide pura). Corrisponde alla percezione di ``pienezza'' del suono.

\range
\begin{itemize}[nosep]
  \item Sinusoide 440 Hz: 10 Hz (energia in una sola bin)
  \item Noise bianco: 2\,888 Hz (distribuzione ampia)
  \item Tanh drive 20: 2\,110 Hz (armoniche sparse)
  \item 2 sinusoidi 200+4000 Hz: 1\,900 Hz (distanza tra le componenti)
\end{itemize}

\oss Lo spread è il descrittore migliore per catturare la distanza
tra componenti spettrali (due sinusoidi lontane) e la ricchezza armonica
della distorsione (tanh drive 1: 174 Hz $\to$ drive 20: 2\,110 Hz).

\rif Peeters (2003, sez.~6.1.2), Lerch (2023, sez.~3.5.2).

% ---
\desc{Rolloff spettrale (\emph{spectral rolloff})}

\formula
La frequenza $f_r$ tale che:
\[
  \sum_{k=0}^{k_r} X(k) = 0{,}85 \cdot \sum_{k=0}^{N-1} X(k)
\]
dove $k_r$ è l'indice della bin corrispondente a $f_r$.

\ascolta
La frequenza sotto la quale si trova l'\textbf{85\% dell'energia}. Indica
il ``confine superiore'' effettivo dello spettro. Complementa il centroide:
il centroide dice dov'è il baricentro, il rolloff dice fin dove arriva
l'energia.

\range
\begin{itemize}[nosep]
  \item Sinusoide 440 Hz: 445 Hz
  \item Noise bianco: 8\,512 Hz (energia fino al limite)
  \item Tanh drive 20: 3\,961 Hz
  \item FM indice 10: 6\,398 Hz
\end{itemize}

\rif Lerch (2023, sez.~3.5.4).

% ---
\desc{Pendenza spettrale (\emph{spectral slope})}

\formula
Regressione lineare della magnitudine sulle frequenze:
\[
  \text{slope} = \frac{\displaystyle\sum_{k=0}^{N-1} \big(f(k) - \bar{f}\big) \cdot \big(X(k) - \bar{X}\big)}
                      {\displaystyle\sum_{k=0}^{N-1} \big(f(k) - \bar{f}\big)^2}
\]
dove $\bar{f}$ e $\bar{X}$ sono le medie delle frequenze e delle magnitudini.

\ascolta
La \textbf{tendenza generale} dello spettro: negativa se l'energia decresce
verso le alte frequenze (tipico degli strumenti acustici), positiva se
cresce. Nei nostri test i valori sono molto piccoli e tutti negativi
(o quasi zero), perché il threshold rimuove le bin deboli.

\rif Lerch (2023, sez.~3.5.6).

% ---
\desc{Deviazione del rapporto fra bande ottavali (\emph{OBSIR-std})}

\formula
\[
  \text{OBSIR}_i = \log_{10} E_{i+1} - \log_{10} E_i
  \qquad
  \text{obsir\_std} = \operatorname{std}(\text{OBSIR})
\]
con $E_i = \sum_{k:\,e_i \le f(k) < e_{i+1}} |X(k)|^2$ ed
edge $e = [200, 400, 800, 1600, 3200, 6400, 10000]$~Hz.

\ascolta
Misura quanto il decadimento spettrale è \textbf{non uniforme} fra le
sei bande ottavali. Valori bassi: lo spettro perde energia con lo
stesso rapporto fra ogni coppia di ottave (decadimento regolare).
Valori alti: alcune bande hanno picchi o buchi marcati rispetto alle
adiacenti.

\range
\begin{itemize}[nosep]
  \item Clarinetto contrabbasso, piano (001): 0.568
  \item Clarinetto contrabbasso, forte (004): 0.672
  \item Timpano, tenuto piano (003): 0.397
  \item Timpano, tenuto forte (006): 0.520
  \item (valori completi sul corpus dopo rigenerazione)
\end{itemize}

\rif Essid, Richard, David (2006), \emph{IEEE TASLP} 14(4).

\bigskip
\noindent\textit{Nota di revisione}: ha sostituito il \emph{spectral
decrease} (Peeters 2004 / Lerch 3.5.5), rimosso dai 16 perché ridondante
con lo slope (Timbre Toolbox 2011, Peeters et al., \emph{JASA} 130(5)).
La scheda d'archivio resta in \texttt{diario/decrease/}.


% ============================================================
\section{Distribuzione}
% ============================================================

Questi descrittori caratterizzano la \emph{forma} della distribuzione
spettrale: quanto è piatta, appuntita, simmetrica.

% ---
\desc{Flatness spettrale (\emph{spectral flatness}, SFM)}

\formula
\[
  \text{SFM} = \frac{\text{media geometrica di } X(k)}
                     {\text{media aritmetica di } X(k)}
  = \frac{\Big(\prod_{k} X(k)\Big)^{1/N}}
         {\frac{1}{N}\sum_{k} X(k)}
\]

In pratica, per evitare underflow:
\[
  \text{SFM} = \frac{\exp\!\Big(\frac{1}{N}\sum_k \ln X(k)\Big)}
                     {\frac{1}{N}\sum_k X(k)}
\]

\ascolta
La \textbf{uniformità} dello spettro. Uno spettro perfettamente piatto
(noise bianco) ha SFM $\approx 1$. Uno spettro con pochi picchi su sfondo
silenzioso (sinusoide) ha SFM basso. Misura la ``rumorosità'' globale.

\range
\begin{itemize}[nosep]
  \item Sinusoide 440 Hz: 0.28
  \item Noise bianco: 0.94
  \item Tanh drive 20: 0.40
  \item FM indice 10: 0.55
  \item Mix sin 50\% + noise 50\%: 0.74
\end{itemize}

\oss La flatness è calcolata solo sulle bin attive (sopra soglia) per
evitare che le bin a zero trascinino il valore. La tanh cambia poco la
flatness (0.28 $\to$ 0.40 da drive 1 a 20) perché le armoniche si
aggiungono come picchi discreti, non come spettro continuo.

\rif Peeters (2003, sez.~9.1), Lerch (2023, sez.~3.5.10), MPEG-7.

% ---
\desc{Crest Factor spettrale (\emph{spectral crest factor})}

\formula
\[
  \text{crest} = \frac{\max_k X(k)}{\frac{1}{N}\sum_k X(k)}
\]
Calcolato solo sulle bin attive (sopra soglia).

\ascolta
Quanto il \textbf{picco massimo} domina sullo spettro.  Alto quando c'è
una componente molto più forte delle altre (fondamentale dominante),
basso quando l'energia è distribuita (noise). È concettualmente
l'inverso della flatness, ma si concentra sul singolo massimo anziché
sulla forma complessiva.

\range
\begin{itemize}[nosep]
  \item Sinusoide 440 Hz: 3.3
  \item Noise bianco: 2.5
  \item Tanh drive 20: 9.4 (fondamentale domina sulle armoniche)
  \item Noise BP Q=50: 12.2 (filtro stretto concentra l'energia)
  \item Mix sin 50\% + noise 50\%: 36.9 (la sinusoide emerge dal noise)
\end{itemize}

\oss Il crest factor cattura bene la distorsione (tanh drive 1: 4.4 $\to$
drive 5: 8.0 $\to$ drive 20: 9.4), perché la fondamentale cresce
rispetto alle armoniche più deboli. Il valore più alto si ha per la
miscela 50/50 (36.9): la sinusoide è un picco netto su un pavimento
di rumore.

\rif Peeters (2003, sez.~9.1), Lerch (2023, sez.~3.5.9).

% ---
\desc{Asimmetria spettrale (\emph{spectral skewness})}

\formula
\[
  \gamma_1 = \frac{\displaystyle\sum_{k=0}^{N-1}
    \left(\frac{f(k) - \mu_1}{\sigma}\right)^3 \cdot X(k)}
    {\displaystyle\sum_{k=0}^{N-1} X(k)}
\]

\ascolta
L'\textbf{asimmetria} dello spettro rispetto al centroide. Positiva se
l'energia è concentrata sotto il centroide (coda verso le alte frequenze),
negativa se l'energia è sopra il centroide (coda verso le basse).

\range
\begin{itemize}[nosep]
  \item Sinusoide 440 Hz: $\sim 0$ (simmetrica, una sola componente)
  \item Noise bianco: $\sim 0$ (simmetrico)
  \item Tanh drive 1: 4.6 (fondamentale forte, poche armoniche deboli in alto)
  \item FM indice 10: $-0.15$ (quasi simmetrica, bande laterali bilaterali)
  \item Mix sin 75\% + noise 25\%: 5.6 (la sinusoide sbilancia verso il basso)
\end{itemize}

\oss La skewness alta della tanh drive 1 (4.6) indica che le poche
armoniche sono molto più deboli della fondamentale e si estendono
verso l'alto. Con drive 20 la skewness scende a 1.7: le armoniche
sono più bilanciate.

\rif Lerch (2023, sez.~3.5.3).

% ---
\desc{Curtosi spettrale (\emph{spectral kurtosis})}

\formula
\[
  \gamma_2 = \frac{\displaystyle\sum_{k=0}^{N-1}
    \left(\frac{f(k) - \mu_1}{\sigma}\right)^4 \cdot X(k)}
    {\displaystyle\sum_{k=0}^{N-1} X(k)} - 3
\]
Il $-3$ (``curtosi in eccesso'') fa sì che una distribuzione
gaussiana dia valore 0.

\ascolta
Quanto lo spettro è \textbf{appuntito} (leptocurtico) o \textbf{piatto}
(platicurtico). Valori alti indicano poche componenti isolate con code
lunghe, valori bassi (negativi) indicano distribuzione uniforme.

\range
\begin{itemize}[nosep]
  \item Sinusoide 440 Hz: 1.5
  \item Noise bianco: $-1.2$ (distribuzione molto piatta)
  \item Tanh drive 1: 19.5 (spettro estremamente appuntito)
  \item Tanh drive 20: 2.2 (meno appuntito, più armoniche)
  \item FM indice 10: $-1.0$ (quasi piatto, tante bande laterali)
\end{itemize}

\oss La kurtosis ha un range enorme (da $-1.2$ a 19.5) e distingue
molto bene il numero di componenti attive. La tanh drive 1
(fondamentale + pochissime armoniche) ha kurtosis altissima.
La FM a indice alto, che ha molte bande laterali di ampiezza simile,
ha kurtosis vicina a $-1$, quasi come il noise.

\rif Lerch (2023, sez.~3.5.3).

% ---
\desc{Entropia spettrale (\emph{spectral entropy})}

\formula
\[
  p(k) = \frac{|X(k)|^2}{\sum_j |X(j)|^2}
\]
\[
  \text{entropy} = \frac{-\sum_k p(k) \log_2 p(k)}{\log_2 N}
\]
con $N$ numero totale di bin (la normalizzazione per $\log_2 N$ porta
il valore in $[0,1]$).

\ascolta
La \textbf{distribuzione di energia} sullo spettro: bassa quando una o
poche bin dominano (sinusoide), alta quando l'energia è distribuita
uniformemente (noise). A differenza di flatness e crest, non satura
sui casi estremi e dilata la zona di mezzo dove vivono i segnali
armonici complessi.

\range
\begin{itemize}[nosep]
  \item Sinusoide bin-esatto: 0.13
  \item Noise bianco: 0.94
  \item Tanh drive 20: 0.22
  \item FM indice 10: 0.55
  \item Mix sin 50\% + noise 50\%: 0.23
  \item Clarinetto contrabbasso piano $\to$ forte: 0.20 $\to$ 0.33
  \item Timpano tenuto: 0.13; attacchi: 0.29--0.32
\end{itemize}

\rif Shen, Hung, Lee (1998); Misra et al.\ (2004), \emph{IEEE ICASSP}.

\bigskip
\noindent\textit{Nota di revisione}: ha sostituito il \emph{coefficiente
di tonalità} in famiglia Distribuzione (riorganizzazione 5+5+3+3 del
17/04). Il tonality coefficient è ora in famiglia Tonalità.


% ============================================================
\section{Tonalità}
% ============================================================

Questi descrittori misurano quanto il segnale è \emph{periodico}
(tonale) rispetto a \emph{casuale} (rumoroso).

% ---
\desc{Rapporto di potenza tonale (\emph{tonal power ratio}, TPR)}

\formula
\begin{enumerate}[nosep]
  \item Si calcola la potenza $P(k) = X(k)^2$
  \item Si trovano i picchi (massimi locali) con \texttt{scipy.signal.find\_peaks},
        parametro \emph{distance} = 3, \emph{prominence} $> 0$
  \item Si filtra per prominenza: solo picchi con prominenza sopra la mediana
  \item Per ogni picco si somma la potenza del lobo principale ($\pm 2$ bin)
  \item $\text{TPR} = \frac{\text{potenza dei picchi}}{\text{potenza totale}}$
\end{enumerate}

\ascolta
La \textbf{frazione di energia} concentrata in picchi spettrali prominenti.
Alto per segnali tonali (tutta l'energia è in picchi), basso per noise
(l'energia è distribuita uniformemente, pochi picchi emergono).

\range
\begin{itemize}[nosep]
  \item Sinusoide 440 Hz: 1.00
  \item Noise bianco: 0.76
  \item Mix sin 50\% + noise 50\%: 0.98
  \item Mix sin 25\% + noise 75\%: 0.90
\end{itemize}

\oss Con l'implementazione basata sulla prominenza, il TPR va da 0.76
(noise) a 1.00 (sinusoide), con un range più ampio della versione
precedente (0.40--0.60). Resta però poco discriminante per segnali
che hanno comunque dei picchi (la tanh ha sempre TPR $\approx 1$).

\rif Lerch (2023, sez.~3.5.11), pyACA (Lerch, implementazione di riferimento).

% ---
\desc{Numero di picchi spettrali}

\formula
Conteggio dei picchi trovati dall'algoritmo TPR (prominenza $>$ mediana,
distanza $\geq 3$ bin).

\ascolta
La \textbf{complessità spettrale}: quanti picchi distinti emergono dallo
spettro. Un segnale con poche armoniche forti ha pochi picchi prominenti,
un noise ne ha molti (casuali).

\range
\begin{itemize}[nosep]
  \item Sinusoide 440 Hz: $\sim$76
  \item Noise bianco: $\sim$96
  \item Tanh drive 1: $\sim$70
  \item Tanh drive 20: $\sim$22
  \item FM indice 10: $\sim$23
  \item 2 sinusoidi 200+4000 Hz: $\sim$27
\end{itemize}

\oss Il conteggio è inversamente correlato alla complessità armonica:
la tanh drive 20 ha \emph{meno} picchi (22) della sinusoide pura (76)
perché le armoniche forti alzano la mediana della prominenza, escludendo
i picchi più deboli. Il noise ha sempre $\sim$96 picchi (il massimo
nella banda 0--10\,000 Hz).

% ---
\desc{Coefficiente di tonalità (\emph{tonality coefficient})}

\formula
\[
  \text{tonality} = 1 - \text{flatness}
\]
con la flatness calcolata come la media geometrica/aritmetica dello
spettro di potenza (vedi sezione Distribuzione).

\ascolta
La \textbf{tonalità} dello spettro: alto quando lo spettro è piccato
(pochi picchi dominanti), basso quando è piatto (rumore). È
complementare alla flatness ma vive in famiglia Tonalità perché viene
usato come indicatore di "quanto suona tonale" più che di forma della
distribuzione.

\range
\begin{itemize}[nosep]
  \item Sinusoide 100 Hz: 0.171
  \item Clarinetto contrabbasso piano (001): 0.152
  \item Clarinetto contrabbasso forte (004): 0.117
  \item Noise bianco: $\approx 0$
\end{itemize}

\rif Peeters (2004, sez.~6.1.7); MPEG-7 Audio (\emph{spectral
flatness}, da cui deriva).

\bigskip
\noindent\textit{Nota di revisione}: questa sezione era originariamente
in famiglia Distribuzione. Spostata in Tonalità con la riorganizzazione
5+5+3+3 del 17/04 (più affine concettualmente a TPR e n. picchi).
Sostituisce la sezione \emph{ACF max} (Lerch 3.5.12), descrittore
rimosso dai 16 al 17/04 perché non separava bene i gesti strumentali
sul corpus. Scheda d'archivio in \texttt{diario/acf/}.


% ============================================================
\section{Dinamica}
% ============================================================

Questi descrittori misurano il \emph{cambiamento} del segnale nel tempo.

% ---
\desc{Flusso spettrale (\emph{spectral flux})}

\formula
\[
  \text{flux}(n) = \sum_{k=0}^{N-1} \big| X_n(k) - X_{n-L}(k) \big|
\]
dove $n$ è il frame corrente e $L$ è il \emph{lag} (default 1 = frame
consecutivi).

\ascolta
La \textbf{quantità di cambiamento spettrale} tra frame. Alto quando lo
spettro cambia rapidamente (noise, glissando veloce), basso per segnali
stazionari (sinusoide costante).

Il parametro \emph{lag} permette di regolare la scala temporale:
con lag~1 ($\sim$43~ms) si cattura il microcambiamento, con lag~8
($\sim$341~ms) si catturano transizioni timbriche più lente.

\range
\begin{itemize}[nosep]
  \item Sinusoide 440 Hz: $\sim$0 (stazionaria)
  \item Noise bianco: 1.30 (cambia continuamente)
  \item Glissando lento: 0.25 (cambiamento moderato)
  \item Glissando veloce: 0.65 (cambiamento rapido)
  \item Tanh drive crescente (10 s): 0.004 (transizione timbrica lenta)
\end{itemize}

\oss Il flux è molto sensibile allo spostamento in frequenza (glissando)
ma poco al cambiamento timbrale graduale (tanh, FM). La tanh drive
crescente ha flux quasi zero perché le armoniche si aggiungono
lentamente e la fondamentale resta ferma. Con lag~8 la FM veloce
raggiunge 0.85, confrontabile col glissando lento (0.84).

\rif Peeters (2003, sez.~6.2.1 ``spectral variation''), Lerch (2023, sez.~3.5.8).

% ---
\desc{Irregolarità spettrale (\emph{spectral irregularity})}

\formula
\[
  \text{irregularity} = \sum_{k=0}^{N-2} \big| \ln X(k) - \ln X(k+1) \big|
\]

\ascolta
La \textbf{frastagliatura} dello spettro: la somma delle differenze
logaritmiche tra bin consecutive. Alta quando lo spettro ha molte
transizioni brusche (noise, con bin alte e basse alternate), bassa
per spettri lisci (sinusoide, poche armoniche isolate).

\range
\begin{itemize}[nosep]
  \item Sinusoide 440 Hz: 29
  \item Noise bianco: 2\,671
  \item Tanh drive 20: 276
  \item FM indice 10: 704
  \item 2 sinusoidi: $\sim$55
\end{itemize}

\oss L'irregolarità è proporzionale al numero di transizioni nel log-spettro.
Due sinusoidi (55) hanno circa il doppio dell'irregolarità di una sola (29),
perché raddoppiano i salti nel log-spettro.

\rif Park (2004), Krimphoff et al.\ (1994).

% ---
\desc{Tasso di attraversamento dello zero (\emph{zero crossing rate}, ZCR)}

\formula
\[
  \text{ZCR} = \frac{1}{N-1} \sum_{n=1}^{N-1}
    \mathbb{1}\!\big[\text{sign}(x(n)) \neq \text{sign}(x(n-1))\big]
\]
dove $x(n)$ è il segnale nel dominio del tempo.

\ascolta
La \textbf{frequenza media approssimativa} del segnale nel dominio del
tempo. Una sinusoide a frequenza $f$ con sample rate $f_s$ ha
$\text{ZCR} = 2f/f_s$. Il noise bianco ha ZCR $\approx 0.5$
(attraversa lo zero circa metà dei campioni).

\range
\begin{itemize}[nosep]
  \item Sinusoide 440 Hz: 0.010 ($= 2 \times 440 / 96000$)
  \item Noise bianco: 0.499
  \item Tanh drive 20: 0.010 (stessa $f_0$, la distorsione non cambia lo ZCR)
  \item FM indice 10: 0.092 (frequenza istantanea varia)
  \item Noise BP Q=50: 0.028
\end{itemize}

\oss Lo ZCR è un descrittore del dominio del tempo, non spettrale.
Discrimina bene tra segnali a bassa e alta frequenza, e tra periodici
e aperiodici. La tanh non cambia lo ZCR perché non sposta la
fondamentale, solo aggiunge armoniche.

\rif Lerch (2023, sez.~3.5.13).


% ============================================================
\section{Tabella riassuntiva}
% ============================================================

La tabella seguente riporta i valori medi di tutti i 16 descrittori
per una selezione di segnali di test (analisi limitata a 0--10\,000 Hz).

\noindent\textit{Nota di revisione (17/04--19/04)}: la tabella è in
attesa di rigenerazione dopo la sostituzione \emph{decrease} →
\emph{OBSIR-std} (Forma) e \emph{ACF max} → \emph{spectral entropy}
(Distribuzione), e lo spostamento del coefficiente di tonalità in
famiglia Tonalità. La sorgente macchina si trova in
\texttt{analisi/tabelle/segnali\_tabelle.md} dopo \texttt{make}.

{
\small
\setlength{\tabcolsep}{3pt}
\begin{longtable}{l|rrrrr|rrrrr|rrr|rrr}
\toprule
& \multicolumn{5}{c|}{\textbf{Forma}} & \multicolumn{5}{c|}{\textbf{Distribuzione}} & \multicolumn{3}{c|}{\textbf{Tonalità}} & \multicolumn{3}{c}{\textbf{Dinamica}} \\
\textbf{Segnale}
& \rotatebox{90}{Centroid}
& \rotatebox{90}{Spread}
& \rotatebox{90}{Rolloff}
& \rotatebox{90}{Slope}
& \rotatebox{90}{OBSIR-std}
& \rotatebox{90}{Flatness}
& \rotatebox{90}{Crest}
& \rotatebox{90}{Skewness}
& \rotatebox{90}{Kurtosis}
& \rotatebox{90}{Entropy}
& \rotatebox{90}{TPR}
& \rotatebox{90}{N.~picchi}
& \rotatebox{90}{Tonality}
& \rotatebox{90}{Flux}
& \rotatebox{90}{Irregularity}
& \rotatebox{90}{ZCR} \\
\midrule
\endfirsthead
\toprule
\textbf{Segnale} & \multicolumn{5}{c|}{\textbf{Forma}} & \multicolumn{5}{c|}{\textbf{Distribuz.}} & \multicolumn{3}{c|}{\textbf{Tonal.}} & \multicolumn{3}{c}{\textbf{Dinam.}} \\
\midrule
\endhead
\bottomrule
\endfoot

Sin 440
& 440 & 10 & 445 & $\sim$0 & 4.3
& .28 & 3.3 & 0.0 & 1.5 & .13
& 1.00 & 76 & .093
& 0.0 & 29 & .010 \\

Noise
& 5000 & 2888 & 8512 & $\sim$0 & .14
& .94 & 2.5 & 0.0 & $-$1.2 & .94
& .76 & 96 & .005
& 1.30 & 2671 & .499 \\

Tanh d.1
& 476 & 174 & 445 & $\sim$0 & 3.4
& .27 & 4.4 & 4.6 & 19.5 & .13
& 1.00 & 70 & .094
& 0.0 & 51 & .010 \\

Tanh d.5
& 918 & 872 & 1336 & $\sim$0 & 2.3
& .29 & 8.0 & 2.3 & 5.5 & .18
& 1.00 & 45 & .090
& 0.0 & 169 & .010 \\

Tanh d.20
& 1876 & 2110 & 3961 & $\sim$0 & 2.0
& .40 & 9.4 & 1.7 & 2.2 & .22
& 1.00 & 22 & .067
& 0.0 & 276 & .010 \\

\midrule

Noise BP 500
& 2486 & 2066 & 4684 & $\sim$0 & .75
& .71 & 7.7 & 1.5 & 1.6 & .73
& 1.00 & 96 & .025
& 1.69 & 2229 & .062 \\

Noise BP 200
& 1803 & 1420 & 3058 & $\sim$0 & .97
& .65 & 9.8 & 2.1 & 4.7 & .62
& 1.00 & 96 & .032
& 1.25 & 1683 & .042 \\

Noise BP 50
& 1288 & 677 & 1788 & $\sim$0 & 1.3
& .57 & 12.2 & 2.4 & 7.4 & .46
& .99 & 96 & .040
& .78 & 924 & .028 \\

\midrule

FM 0.5
& 545 & 323 & 1090 & $\sim$0 & 3.0
& .35 & 5.9 & 1.7 & 2.5 & .20
& 1.00 & 77 & .075
& 0.0 & 125 & .009 \\

FM 3
& 1422 & 864 & 2414 & $\sim$0 & 2.8
& .44 & 4.3 & 0.6 & $-$0.1 & .42
& 1.00 & 58 & .060
& 0.0 & 322 & .028 \\

FM 10
& 4283 & 2238 & 6398 & $\sim$0 & .75
& .55 & 3.6 & $-$0.2 & $-$1.0 & .55
& 1.00 & 23 & .043
& 0.0 & 704 & .092 \\

\midrule

Sin+N 75/25
& 701 & 1231 & 453 & $\sim$0 & 2.1
& .20 & 8.1 & 5.6 & 33.8 & .14
& 1.00 & 96 & .117
& .04 & 237 & .057 \\

Sin+N 50/50
& 3671 & 3193 & 7884 & $\sim$0 & 1.6
& .74 & 36.9 & 0.5 & $-$1.2 & .23
& .98 & 96 & .022
& .70 & 2703 & .179 \\

Sin+N 25/75
& 4631 & 3027 & 8350 & $\sim$0 & .98
& .90 & 19.0 & 0.1 & $-$1.3 & .58
& .90 & 96 & .008
& 1.08 & 3176 & .454 \\

\midrule

Gliss lento
& 1098 & 11 & 1109 & $\sim$0 & 5.5
& .33 & 3.2 & 0.0 & 0.9 & .14
& 1.00 & 89 & .080
& .25 & 29 & .023 \\

Gliss vel.
& 1096 & 16 & 1111 & $\sim$0 & 5.3
& .39 & 3.3 & 0.0 & 0.7 & .18
& .99 & 86 & .067
& .65 & 29 & .023 \\

\midrule

2 sin 200+4k
& 2131 & 1900 & 4008 & $\sim$0 & 8.4
& .42 & 2.4 & 0.0 & $-$2.0 & .23
& 1.00 & 27 & .062
& 0.0 & 55 & .091 \\

2 sin 400+1k
& 703 & 300 & 1008 & $\sim$0 & 3.9
& .40 & 2.6 & 0.0 & $-$2.0 & .23
& 1.00 & 26 & .066
& 0.0 & 55 & .029 \\

\midrule

Tanh cresc
& 1346 & 1446 & 2678 & $\sim$0 & 2.1
& .34 & 8.7 & 2.1 & 4.5 & .20
& 1.00 & 55 & .079
& .004 & 226 & .010 \\

FM cresc
& 2339 & 1303 & 3704 & $\sim$0 & 1.9
& .50 & 3.9 & 0.3 & $-$0.4 & .45
& 1.00 & 53 & .051
& .04 & 443 & .046 \\

\end{longtable}
}


% ============================================================
\section{Quale descrittore per cosa?}
% ============================================================

\subsection{Riconoscere la brillantezza}
Il \textbf{centroide} è il descrittore principale. Il \textbf{rolloff}
lo complementa indicando il confine superiore dello spettro.

\subsection{Riconoscere la ricchezza armonica}
Lo \textbf{spread} è il più sensibile: segue la distorsione
(tanh drive 1: 174 Hz $\to$ drive 20: 2\,110 Hz) e la distanza
tra componenti. La \textbf{kurtosis} complementa con la forma:
alta per poche armoniche (19.5), bassa per molte ($-1.0$).

\subsection{Distinguere tonale da rumoroso}
La \textbf{spectral entropy} dà la misura più graduata e con range
ampio (0.13 sinusoide $\to$ 0.94 noise), senza saturare nella zona di
mezzo dove vivono i segnali armonici complessi (FM e tanh fra 0.2 e
0.6). La \textbf{flatness} è complementare ma satura sui casi estremi
(0.945 sul bin esatto, 0.99 sul noise). Il \textbf{coefficiente di
tonalità} (in famiglia Tonalità) è $1 - \text{flatness}$ riscritta in
chiave tonale. Il \textbf{TPR} discrimina bene sui segnali a picchi
prominenti, lo \textbf{ZCR} conferma nel dominio del tempo.

\subsection{Catturare il cambiamento}
Il \textbf{flux} misura il cambiamento spettrale tra frame, con la
possibilità di regolare la scala temporale tramite il parametro lag.
Lo \textbf{ZCR} cattura cambiamenti nella frequenza fondamentale.

\subsection{Leggere il decadimento spettrale}
Lo \textbf{slope} dà la pendenza media (informazione globale, può
essere rumorosa frame-per-frame). L'\textbf{OBSIR-std} misura quanto
il decadimento è non uniforme fra ottave: complementare allo slope,
sale quando alcune bande "saltano" rispetto alle adiacenti (risonanze
o buchi).

\subsection{Descrittori archiviati}
Lo \textbf{spectral decrease} (Peeters 2004 / Lerch 3.5.5) e l'\textbf{ACF
max} (Lerch 3.5.12) sono stati rimossi dai 16: il primo per ridondanza
con lo slope (Timbre Toolbox 2011), il secondo perché non separava
bene i gesti strumentali sul corpus. Le schede d'archivio restano in
\texttt{diario/decrease/} e \texttt{diario/acf/} come riferimento storico.

\end{document}
